\documentclass[11pt,a4paper]{article}
\usepackage[UTF8]{ctex}
\usepackage{amsmath,amssymb,amsthm}
\usepackage{geometry}
\usepackage{enumitem}
\usepackage{xcolor}
\usepackage{tcolorbox}
\usepackage{fancyhdr}
\usepackage{titlesec}
\usepackage{hyperref}

\geometry{margin=1in}
\setCJKmainfont{SimSun}[AutoFakeBold=true]

\pagestyle{fancy}
\fancyhf{}
\fancyfoot[C]{CS 70, 2026年春季, 作业 00 解答}
\renewcommand{\headrulewidth}{0pt}

\titleformat{\section}{\large\bfseries}{\thesection}{1em}{}

\definecolor{solutiongreen}{RGB}{220,252,231}
\definecolor{exampleblue}{RGB}{248,249,250}

\newtcolorbox{solutionbox}{
    colback=solutiongreen,
    colframe=green!50!black,
    title=解答,
    fonttitle=\bfseries\small
}

\newtcolorbox{examplebox}{
    colback=exampleblue,
    colframe=gray!30,
    boxrule=0.5pt
}

\begin{document}

\begin{center}
    {\Large\bfseries CS 70}\\[0.3em]
    {\large 离散数学与概率论}\\[0.2em]
    2026年春季\\[0.2em]
    Sinclair, Song\\[1em]
    \fcolorbox{green!50!black}{green!10}{\parbox{0.4\textwidth}{\centering
        \textbf{作业 00 解答}
    }}
\end{center}

\section{行政事务}

\begin{enumerate}[label=(\alph*)]
    \item 课程网站位于 \url{https://www.eecs70.org/}
    
    \item 
    \begin{solutionbox}
        \begin{enumerate}[label=(\roman*)]
            \item 讨论出勤：5\%，迷你维生素：5\%，作业：15\%，期中考试：30\%，期末考试：45\%。
            \item 本学期 CS70 没有免作业选项。
            \item 每次讨论班你将获得 1 个出勤点，需要至少 13 个点才能获得讨论出勤的满分。欢迎参加其他讨论班，任何讨论班都可以获得学分，但强烈建议你持续参加学期初注册的讨论班。
            \item 当周的作业在周一之前发布在课程网站上。作业需在周六 4:00 PM 前提交到 Gradescope（宽限期至 6:00 PM）；该作业的解答将在周一发布。
            \item 整个学期你可以丢弃最低的 3 次作业，最高的 13 次迷你维生素将计入成绩。但请将这些丢弃机会留给紧急情况。
            \item 期中考试日期：03/12/26 周四 7--9pm\\
                  期末考试日期：05/14/26 周四 3--6pm
            \item 73\%。
        \end{enumerate}
    \end{solutionbox}
\end{enumerate}

\section{课程政策}

\begin{enumerate}[label=(\alph*)]
    \item \textbf{是}，这违反了课程政策。所有解答必须完全由提交作业的学生撰写。即使学生合作，每个学生也必须撰写独特、独立的解答。在这种情况下，Alice 和 Bob 都将承担责任。
    
    \item \textbf{否}，这不违反课程政策。虽然分享书面解答是不允许的，但分享问题的方法是允许并鼓励的。因为 Carol 只是抄了笔记，而不是 Dan 的解答，并且正确注明了 Dan 的贡献，这是一种被积极鼓励的合作形式。
    
    \item \textbf{否}，这不违反课程政策。使用外部资源帮助解决作业问题，只要：(1) 学生确保理解解答；(2) 学生利用理解写出新的解答，而不是从外部来源复制；(3) 学生注明外部来源，这是可以的。然而，在网上查找作业问题是违反课程政策的；在网上找到作业解答后，正确的做法是关闭标签页。
    
    \item \textbf{是}，这违反了课程政策，Frank 和 Grace 都将承担责任。尽管 Frank 注明了 Grace，但书面解答根本不应该被分享，更不应该被抄下来。这是为了确保每个学生学会如何撰写和呈现清晰、有说服力的论证。
    
    \item \textbf{是}，这违反了课程政策。再次强调，引用并不能弥补书面解答不应该被分享的事实，无论是书面还是打字形式。在这种情况下，Heidi 和 Irene 都将承担责任。
    
    \item \textbf{是}，这违反了课程政策。Joe 不应该在解答正式发布之前阅读它们。相反，当 Joe 卡住时，应该通过 Ed 或办公时间寻求帮助。
    
    \item \textbf{是}，这违反了课程政策。Kai 不应该直接让 ChatGPT 解决作业问题。如果 Kai 真的想用 LLM 来帮助做作业，他们应该问一个概念性问题（他们会在办公时间或 Ed 上合理提出的问题）。
\end{enumerate}

\section{Ed 的使用}

\begin{enumerate}[label=(\alph*)]
    \item \begin{solutionbox}
        这个问题有两个问题：
        \begin{enumerate}[label=\arabic*.]
            \item 这个问题没有通过 5 分钟测试。这个概念需要超过 5 分钟来解释，因此更适合在办公时间讨论。
            \item 这个问题没有聚焦于学生具体困惑的概念。Ed 上的问题应该具体，并应包括导致问题出现的每一步推理。一个更好的问题可能是："我理解了笔记 2 中定理 XYZ 证明的每一步，除了最后一步。我尝试这样推理，但我不明白它如何得出结果。有人能解释我哪里出错了吗？"
        \end{enumerate}
    \end{solutionbox}
    
    \item 每周公告在每周一发布。它们是必读内容。
    
    \item 请发送电子邮件至 \texttt{sp26@eecs70.org}。
\end{enumerate}

\section{学术诚信}

（此题为签名题，无需解答）

\section{命题逻辑练习}

\textcolor{blue}{\textbf{笔记 1}}

\begin{enumerate}[label=(\alph*)]
    \item \begin{solutionbox}
        这道题最棘手的部分是"唯一"这个词。我们可以用两个用"且"连接的陈述来表达唯一解的存在：(1) 存在一个解，(2) 存在的任意两个解必须相同。
        
        因此，我们可以将这个陈述重写为："对于每个实数 $k$，存在一个实数 $x$ 使得 $x^3 = k$，并且对于所有实数 $y$ 和 $z$，如果 $y^3 = k$ 且 $z^3 = k$，则 $y = z$。"
        
        用命题逻辑表示为：
        \[
        (\forall k \in \mathbb{R})\left[(\exists x \in \mathbb{R})(x^3 = k) \land (\forall y,z \in \mathbb{R})[((y^3 = k) \land (z^3 = k)) \Rightarrow (y = z)]\right]
        \]
    \end{solutionbox}
    
    \item \begin{solutionbox}
        这句话可以用命题逻辑写成：
        \[
        (\forall p \in \mathbb{N})[P(p) \Rightarrow ((\forall a,b \in \mathbb{N})[((p \mid ab) \land \neg(p \mid a)) \Rightarrow (p \mid b)])]
        \]
    \end{solutionbox}
    
    \item \begin{solutionbox}
        如果两个实数的乘积为 0，那么其中必有一个为 0。
    \end{solutionbox}
    
    \item \begin{solutionbox}
        不存在一个自然数能整除所有大于它的合数。
    \end{solutionbox}
\end{enumerate}

\end{document}
